\section{Introduction}
\label{sec:intro}

Demand-responsive transit (DRT) has emerged as a promising solution for providing
flexible, on-demand mobility in low-density urban areas where fixed-route bus services
are economically unviable.
In China, rapid suburbanisation has created large residential zones with sparse transit
coverage, making DRT an increasingly important policy instrument for urban planners
and transport authorities.
However, the dominant door-to-door DRT model — in which vehicles travel directly
between each passenger's origin and destination — is operationally costly: vehicles
accumulate large detour distances serving spatially dispersed requests, and fleet
utilisation remains low.
Meeting-point-based DRT addresses this inefficiency by requiring passengers to walk
to and from pre-defined consolidation points (e.g., bus stops, commercial landmarks),
enabling vehicles to serve multiple passengers along shared route segments.
The trade-off is clear: meeting points reduce vehicle-kilometres but impose a walking
burden on passengers, and passengers who find the required walk unacceptable will
reject the service offer and use an alternative mode.
Jointly optimising the spatial assignment of meeting points and the vehicle routing
schedule — while explicitly accounting for passenger heterogeneity in mode choice —
is therefore the central operational challenge for next-generation DRT systems.

Existing literature on meeting-point DRT has made important progress but leaves a
critical gap.
The Dial-a-Ride Problem (DARP) literature provides a rich foundation for vehicle
routing with time windows and capacity constraints
\citep{cordeau2007,ho2018}.
Recent extensions have introduced meeting points into the DARP framework.
\citet{cortenbach2024} proposed the DARPmp, the first formulation of DARP with
meeting points, using a Tabu Search heuristic to solve instances with single-sided
meeting point assignment: pickup locations are flexible (passengers walk to a nearby
stop), but dropoff locations remain fixed at the passenger's destination.
This single-sided approach captures part of the efficiency gain but leaves the dropoff
leg unoptimised.
Crucially, \citet{cortenbach2024} do not model passenger choice: all offered meeting
points are assumed to be accepted, ignoring the reality that passengers with low
walking tolerance will reject inconvenient offers.
On the dynamic side, \citet{wu2025} proposed a rolling horizon framework for
online DRT re-optimisation, demonstrating that periodic re-assignment of requests
to vehicles significantly reduces operational cost under dynamic demand.
However, \citet{wu2025} do not incorporate meeting points: their model retains
door-to-door service.
No existing work combines: (a)~bidirectional meeting point sets ($M_r^P$ for pickup,
$M_r^D$ for dropoff), (b)~explicit binary logit passenger choice with
heterogeneous passenger types and an outside option, and (c)~rolling horizon
re-optimisation in a many-to-many DRT setting.
This is the gap the present paper fills.

We propose a many-to-many DRT model with bidirectional meeting point sets.
For each request $r$, the system constructs a candidate pickup set
$M_r^P \subseteq \mathcal{M}$ (meeting points within walking radius $\rho^P$ of the
passenger's origin) and a candidate dropoff set
$M_r^D \subseteq \mathcal{M}$ (within $\rho^D$ of the destination).
The system jointly optimises the selection of pickup meeting point $m_r^P \in M_r^P$
and dropoff meeting point $m_r^D \in M_r^D$ for each accepted request, together with
vehicle routing and scheduling.
Passenger response is modelled via a binary logit model (special case of MNL) with three
passenger types — price-sensitive, time-sensitive, and walk-sensitive — each with
distinct utility coefficients over walking time, waiting time, in-vehicle travel time,
and fare.
An outside option allows passengers to reject all offered bundles and use an
alternative mode.
The full model is formulated as a coupled three-layer system:
(1)~service offer generation, which enumerates feasible bundles
$\mathcal{B}_r = M_r^P \times M_r^D$ and computes binary logit choice probabilities;
(2)~passenger response, which determines the accepted bundle based on the binary logit acceptance rule;
and (3)~dynamic dispatch, which routes vehicles using a rolling horizon ALNS
heuristic with re-optimisation at fixed time intervals.
We solve small instances exactly using a Mixed-Integer Linear Program (MILP)
implemented in Gurobi, and large instances with the rolling horizon ALNS heuristic.
Numerical experiments on synthetic scenarios calibrated to Chinese suburban conditions
validate the model and quantify the efficiency and equity implications of bidirectional
meeting point assignment.

The contributions of this paper are as follows:

\begin{enumerate}
  \item First formulation of many-to-many DRT with bidirectional meeting point sets
        ($M_r^P$, $M_r^D$) and binary logit passenger choice as a coupled three-layer model.
  \item MILP ex-post routing diagnostic and rolling horizon ALNS heuristic for the bidirectional
        DARP with passenger choice.
  \item Empirical demonstration of a \textbf{35.0\%} vehicle efficiency gain in
        per-trip vehicle utilisation at equal served share (FullModel 11.1 vs.\
        DoorToDoor\,(capped) 17.1\,vkm/trip, endogenous matched-coverage comparison),
        and a 29.2\% gain under the unconstrained comparison (15.1 vs.\ 21.3\,vkm/trip).
  \item Equity analysis across three passenger types (Gini coefficient = 0.1216) and
        five actionable policy recommendations for Chinese city DRT operators.
\end{enumerate}

The remainder of this paper is organised as follows.
Section~\ref{sec:literature} reviews related work on DARP, meeting-point DRT, and
passenger choice modelling.
Section~\ref{sec:model} presents the problem formulation, including the three-layer
coupled model and the binary logit choice structure.
Section~\ref{sec:algorithm} describes the solution methodology: the ex-post MILP routing
diagnostic and the rolling horizon ALNS heuristic.
Section~\ref{sec:experiments} reports numerical experiments, including ablation
studies and sensitivity analysis across demand density, walking tolerance, and fleet
size.
Section~\ref{sec:policy} discusses policy implications for Chinese city DRT operators.
Section~\ref{sec:conclusion} concludes and outlines directions for future research.
